\section{Introduction}

\subsection{Problem Statement and Motivation}







The evergrowing demand for more data generation across all sectors, from industry to healthcare,
has created a need for fast and reliable ways to interpret and understand all of the information 
being generated. Machine learning (ML) has been developed as a solution to address this need. It 
is able to learn patterns in complex and convoluted data without being explicitly programmed for
the specific tasks. This ability makes ML very versatile and capable of being deployed in drastically
different enviroments. Traditional ML models like Support Vector Machines (SVM) or decision trees,
still require significant manual work of cleaning the data, extracting features and fine-tuning parameters.
In recognition of these shortcomings a new type of ML, deep learning (DL) has evolved.
At its core DL models have the ability to automatically derive relevant features from raw data.
The multi-layered approach leverages a series of connected transformations where each layer acts as a filter, 
iteratively refining the input data. This allows the network to build a hierarchy of knowledge, 
where simpler features are combined into intricate concepts, 
ultimately generating a robust and highly abstract representation optimal for prediction or classification.



Machine learning (ML) is a very powerful tool because of its ability to understand and interpret
large sets of data fast and efficiently, to learn patterns without it being explicitly programmed.
While classical ML models learn directly on the provided data,
deep learning is able to learn and understand patterns on multiple layers of abstraction,
called Neural networks with multiple layers (hence the term "deep").
Deep learning has several benefits over traditional machine learning models,
in particular it excels at processing raw input data and it is able
to understand minute as well as broader patterns.
Furthermore it 
These benefits are particularly useful when trying to analyze biosignals to
identify specific states of the human body.

Epilepsy detction and prediction is a field where deep learning could
significantly improve research 